\section{Introduction}

Natural-language video retrieval connects a user's description with relevant content in a video collection. A query may specify an object's appearance, a human action, or a broader scene, requiring retrieval systems to preserve different aspects of visual evidence. Large-scale image--text pretraining provides a foundation for matching text to visual content \cite{radford2021clip}, and subsequent work has investigated how image--language representations can be transferred to video retrieval \cite{luo2021clip4clip}. Nevertheless, a single representation or scoring pathway may not capture every detail expressed in a query.

Automatically generated captions provide an additional retrieval pathway. Instead of comparing a query only with visual features, a system can also compare it with textual descriptions generated from sampled video frames. This pathway may express activities or relations in a form that is easier to match linguistically. However, it also transforms the original evidence: a caption may omit a small object, replace a precise attribute with a general description, or introduce unsupported content. The usefulness of caption evidence therefore needs to be measured relative to the visual branch rather than assumed from the presence of an additional modality.

This study extends our earlier multimodal video retrieval system [PriorWork: bibliographic details pending] by investigating the fusion policy applied to visual and generated-caption retrieval. A fixed policy assigns the same relative importance to the two branches for every query. Such a policy is simple and may generalize well when supervision is limited, but it cannot explicitly respond to differences among user requests. A query-dependent policy can instead predict fusion weights from the query representation. Its additional flexibility raises a corresponding question: does it learn a transferable preference between evidence sources, or merely fit the particular queries used for tuning?

We study this question with frozen retrieval encoders and controlled fusion comparisons. The initial learned design uses a lightweight multilayer perceptron (MLP) to map a query embedding to visual and caption weights. It is compared with an interpretable semantic-category routing policy, a linear gate, and fixed rank-fusion baselines. Reciprocal Rank Fusion (RRF) \cite{cormack2009rrf} provides a common aggregation framework when the underlying branch scores are not directly calibrated. A globally optimized fixed weight is a central control, because outperforming equal weighting alone would not establish the value of query adaptation.

The study addresses three research questions. \textbf{RQ1:} How does the contribution of generated-caption retrieval vary across queries and semantic categories? \textbf{RQ2:} Can a lightweight query-conditioned gate improve held-out retrieval over globally tuned fixed fusion? \textbf{RQ3:} What additional benefit, if any, does a nonlinear MLP provide over semantic routing and linear gating under the same supervision and retrieval inputs?

The intended contributions are a controlled evaluation of query-conditioned fusion for frozen visual and generated-caption branches; a comparison of interpretable, linear, and nonlinear weighting policies; and an analysis of effectiveness, failure cases, and online cost. Query-conditioned expert weighting already appears in video retrieval \cite{miech2018mee}, so our contribution is not the introduction of that general mechanism. Rather, we investigate its behavior in the specific setting of generated-caption augmentation, including whether its gains persist against strong fixed controls and under limited supervision. Conclusions about effectiveness will be drawn from the experiments rather than assumed in the method design.

